\section{Obstructions to retaining the limiting module}\label{sec:obstructions}
\subsection{The profinite comparison and bounded occurrence}
The rational phase channels of Theorem~\ref{thm:packets} define
the auxiliary limiting Hilbert space
\begin{equation}
\cK_\rho=\bigoplus_{\alpha\in2\pi\Q/(2\pi\Z)}L^2(\R),
\qquad
\mathcal V_a^\rho(e_\alpha f)=a^{-1/2}
\sum_{a\beta=\alpha}\sqrt{\rho(\beta)/\rho(\alpha)}
\,e_\beta U_{\log a}f .
\end{equation}
It records limiting mixed data; it is not declared a new physical
sector. Let $D_\rho e_\alpha f=\sqrt{\rho(\alpha)}e_\alpha f$.
Then $\mathcal V_a^\rho=D_\rho\mathcal V_a^1D_\rho^{-1}$.
Fourier series on the profinite integers identify the Haar version
with $L^2(\R\times\widehat\Z)$ and
\begin{equation}\label{eq:profinite-action}
(\mathcal V_a^1F)(x,z)=
\sqrt a\,\mathbf1_{a\widehat\Z}(z)F(x-\log a,z/a).
\end{equation}
It is an isometry, with range projection
$\mathbf1_{a\widehat\Z}$. The comparison makes a precise occurrence
test possible while leaving the YM state fixed.

\begin{theorem}[No bounded class-sector intertwiner]\label{thm:bounded-occurrence}
Every bounded map $T:\cK_\rho\to\cH_{\rm cl}$ satisfying
$T\mathcal V_a^\rho=V_aT$ for all a is zero.
The conclusion also holds on the closed identity-phase cyclic
submodule generated by all $\mathcal V_a^\rho e_0f$.
\end{theorem}
\begin{proof}
The coefficient map $\mathcal C:\ell^2(\N)\to\cH_{\rm cl}$,
$e_n\mapsto\chi_n$, is a bounded isomorphism by
\eqref{eq:gram}. It intertwines $V_a$ with the shift
$S_ae_n=e_{an}$. Thus $A=\mathcal C^{-1}TD_\rho$ obeys
$A\mathcal V_a^1=S_aA$. For $\eta_n=A^*e_n$ and every prime
$p\nmid n$, $(\mathcal V_p^1)^*\eta_n=A^*S_p^*e_n=0$.
By the range projection in \eqref{eq:profinite-action},
$\eta_n$ vanishes on $p\widehat\Z$ for every such p.
The complement of these divisibility sets has Haar measure at
most $\prod_{p\in E}(1-1/p)$ for each finite set E of primes
not dividing n. These products tend to zero by Euler's
divergence of $\sum_p1/p$. Hence $\eta_n=0$ and A=0.

For the cyclic submodule, Haar coordinates identify it with the
closed span of functions $f(x)\mathbf1_{a\widehat\Z}(z)$.
Multiplication by $\mathbf1_{p\widehat\Z}$ preserves this subspace
because intersections are $\operatorname{lcm}(a,p)\widehat\Z$.
The same isometry adjoint and range-projection argument applies
there. Tests in $\cD^0$ are dense in $L^2(\R)$: a vector
orthogonal to $\T\cD$ is an \Ltwo{} distributional solution of
$\T u=0$, hence zero. They generate the same cyclic closure.
\end{proof}
In particular, a uniformly bounded family of recovery maps
cannot have nonzero strong limits on the whole limiting module
while asymptotically preserving all winding relations: its
strong limit would be the forbidden bounded intertwiner.
The theorem concerns a bounded map into the actual \emph{class}
sector with these specified relations. It is not a no-go theorem
for every possible YM source, or for a different physical
representation.

\subsection{What a fixed closable operation can preserve}
\begin{lemma}\label{lem:closable-weak}
Let T be a fixed densely defined closable operator between Hilbert
spaces. If $u_n\in\Dom T$, $u_n\rightharpoonup0$, and
$\sup_n\norm{Tu_n}<\infty$, then $Tu_n\rightharpoonup0$.
\end{lemma}
\begin{proof}
For every $v\in\Dom T^*$,
$\ip v{Tu_n}=\ip{T^*v}{u_n}\to0$.
The adjoint domain is dense by closability; boundedness extends
the limit to every v.
\end{proof}
Thus a fixed closable insertion cannot recover a nonzero strong
vector from the raw escaping packets while keeping their output
norms bounded. This does not forbid a scale-dependent unitary
recentring, which will succeed for the ordinary \Ltwo{} pairing.

There is a second, global obstruction even after that recovery.
\begin{theorem}[No closable global \Ltwo{} completion of the Weil norm]\label{thm:nonclosable}
Let $J_0:L^2(\R)\to\cH$ be any isometry. There is no fixed
closable operator T with $J_0\cD^0\subset\Dom T$ such that
\begin{equation}
\ip{TJ_0f}{TJ_0g}=Q(f,g)\qquad(f,g\in\cD^0).
\end{equation}
\end{theorem}
\begin{proof}
The assumed identity first implies positivity and hence RH by
Weil's criterion. This is a consequence in a contradiction proof,
not an extra hypothesis. Under that consequence the explicit
formula reads
\begin{equation}\label{eq:zero-pairing}
Q(f,g)=\sum_{\gamma\ {\rm distinct}}m_\gamma
             \overline{\widehat f(\gamma)}\widehat g(\gamma).
\end{equation}
Fix one ordinate $\gamma_0$ and $h\in\cD$ with $\widehat h(0)=1$.
Set
\begin{equation}
f_L=\frac{\T}{\gamma_0^2+1/4}
        \left[L^{-1}e^{\ii\gamma_0x}h(x/L)\right].
\end{equation}
Then $\norm{f_L}_2\to0$, while
\[
\widehat f_L(\gamma)=
\frac{\gamma^2+1/4}{\gamma_0^2+1/4}
\,\widehat h(L(\gamma-\gamma_0))
\]
converges in $\ell^2(\{m_\gamma\})$ to the unit point vector
at $\gamma_0$. To justify the norm convergence, use isolation
of $\gamma_0$, the Schwartz estimates of $\widehat h$, and
$N(T)=O(T\log T)$; a fixed summable majorant controls all
other ordinates for $L\ge1$. Hence $TJ_0f_L$ is Cauchy with
nonzero limiting squared norm $m_{\gamma_0}$, even though
$J_0f_L\to0$. This contradicts closability.
\end{proof}
The support of $f_L$ grows. The conclusion concerns closability
relative to the ordinary global \Ltwo{} input norm; it does not
preclude continuous distributional sources on fixed-support
smooth-test spaces, or closable physical operators in a
stronger source topology.
